\section{Introduction}
A finite command catalog creates a distinction between the resources allocated to a decision and the number of commands used to implement it. A service provider may deliver different expected obligations to different customer histories while installing only a small common set of terminal commands. Installing a command incurs a fixed charge. Each history also has an expected cap and a ceiling on every realized obligation. Optimizing the expected targets without optimizing their shared implementation can therefore miss both a discrete cost and a feasibility restriction.

We study this joint allocation and implementation problem, preserving the accepted aggregate promise and deterministic service chosen before the terminal draw. Our contribution is a structural theory rather than an empirical claim about a particular provider. The service-credit interpretation makes the timing and resource accounting concrete. It does not supply customer-choice, incentive-compatibility, or field-calibration assumptions that are absent from the model.

The analysis has three connected parts. First, we eliminate the history-level lotteries constructively and obtain an exact resource allocation problem on a selected path. Each edge has a capacity that is the difference of two values of a common potential. We characterize this conservation property on arbitrary acyclic graphs, give a linear-time recognition algorithm with a path witness when it fails, and identify exactly when complementing resources permits all starting vertices to share one terminal deficit. Thus the useful abstraction is not a renewal-specific identity alone: it is a recognizable class of endogenous path and resource-allocation problems.

Second, we establish a limited-menu complexity frontier inside the original model. Exact optimization is hard in the parameterized sense even when both the command allowance and the number of distinct expected caps are parameters. The reduction retains a common linear reward, free service, uniform probabilities, equal expected caps and realization ceilings, and nonnegative opening charges. Its dedicated baseline histories force the baseline commands to consume the actual budget. This closes the gap between a nonbinding-budget subset-sum construction and the previously established small-type enumeration bound. The lower bound does not contradict polynomial algorithms for a fixed parameter value or numerical-lattice algorithms: it distinguishes their dependence on parameters and encoding.

Third, resource conservation supports an original-feasible additive approximation. A single deficit coordinate merges all feasible starting points. Rounding and repairing on the same selected path preserves the exact promise and the command count. Centered concave edge kernels permit fast convolution even when intermediate states contain unreachable gaps. The abstract theorem covers conservative acyclic paths; the renewal model supplies an exact implementation of its edge allocations. Optimal prototype selection in a declared reward family makes the type approximation operational. Joint reward types, exact lattice regimes and price-based certificates provide complementary refinements in the electronic companion.

The computational analysis is organized around what these results establish: the cost of representing, solving, and checking a decision, not universal superiority over mixed-integer optimization. We preserve all earlier failures and timings and distinguish them from new executions. A direct mixed-integer convex formulation has no service-envelope approximation. Its numerical solver bounds are reported separately from independently checked rational intervals. Screening is treated as an optional preprocessing decision whose cost must be included; its unfavorable earlier matched comparison is retained. The article consequently makes a theory claim supported by algorithmic evidence, not an unverified claim of operational adoption.

\subsection{Relationship to prior work}
Constrained-path methods provide the algorithmic context. Lagrangian bounds and constrained shortest-path search are established tools \citep{HandlerZang1980,BeasleyChristofides1989,Fisher1981}. Approximation schemes for restricted shortest paths also precede this work \citep{Hassin1992,LorenzRaz2001}. Those results do not themselves eliminate our individual expected and realized constraints. Once an exact path representation has been obtained, however, resource-state dynamic programming and Lagrangian branching are antecedent machinery, not new inventions. Our contribution at that interface is the implementable edge model, a testable source-conservation criterion, and a repair that keeps the selected implementation and original equality.

A fixed path leads to separable nonlinear resource allocation; its continuous allocation step belongs to the literature surveyed by \citet{Patriksson2008}. Concave convolution uses matrix-searching structure \citep{Aggarwal1987}; our finite-support argument explains why unreachable holes do not invalidate that use. The new parameterized reduction, rather than the existence of a resource dynamic program, establishes the discrete frontier. The relevant distinction between a parameter-dependent exponent and fixed-parameter tractability follows \citet{DowneyFellows1995b}.

Small support and a small shared catalog are different objects. Extreme-point results for moment-constrained probability measures bound the support of a distribution under specified moment constraints \citep{Winkler1988}. They do not bound the union of history-dependent supports when the union incurs opening charges. Auction menu complexity instead counts outcome--payment alternatives subject to incentive constraints \citep{HartNisan2013,Babaioff2016}. Our terminal commands carry neither a free history label nor a payment menu; histories are observed by the allocator, not selected strategically by a buyer. These distinctions prevent a support bound or an auction approximation theorem from being imported as a proof of our cap-count result.

There is also a distinction from assortment and data-driven optimization. In choice-based assortment optimization, offering an item changes demand through the choice model, and empirical demand learning may be part of the decision problem \citep{Rusmevichientong2010}. Here the history probabilities and accepted promise are model inputs and command probabilities are controlled decisions. Distributionally robust optimization addresses uncertainty in distributions and provides statistical or worst-distribution guarantees under explicit ambiguity models \citep{EsfahaniKuhn2018}; our reward-oscillation certificate holds feasibility inputs fixed and is not such a guarantee. Finally, prototype selection is a clustering problem. We identify its precise oscillation distance and cap restrictions, use the standard facility-location formulation for general medoids, and solve a one-dimensional common-curvature case by weighted-median dynamic programming. We do not claim a new general metric clustering algorithm \citep{Arya2004,Wu1991}.

Table~\ref{tab:antecedents59} states the nearest theorem-level comparisons. Stripping away every certificate, build manifest, and repository artifact leaves the constructive representation, the original-model parameterized hardness theorem, the conservative-path recognition and approximation results, and the loss-controlled implementation frontier. Reproducibility supports these claims but is not a substitute for them.
